\documentclass[thesis]{cluu}

\usepackage[style=cluu]{biblatex}
\addbibresource{pal.bib}

\begin{document}
\title{Analysis and Detection of LLM-Generated Text in Swedish}
\author{Gina Welsh}
% \date{...}
\supervisor{Birger Moëll, Uppsala University}
% Use \supervisor (without s at end) if you have only one

\maketitle

\begin{abstract}

This thesis investigates investigate systematic differences between text generated
by large language models (LLMs) and human-written text in the Swedish language. In addition, the
thesis analyses significant differences in the ability to detect LLM-generated
Swedish text between humans and AI-detection tools. The study will include both a formal and informal text
component, that is, linguistic analyses and AI-detection experiments will be performed on two sources
of data from formal and informal domains.
\end{abstract}

\tableofcontents

\addchap{Preface}
I would like to thank the following people...
% with \addchap instead of \chapter it isn't numbered.

\chapter{Introduction}
The invention of transfomer-based Large Language Models (LLMs) has created a new wave of textual data that is entirely generated by machines and very closely emulates, to the naked eye, text written by humans. As these tools continue to generate unbelievable amounts of textual data at a break-neck speed (<< insert some figure taken recently on number of TB per year >>), changing the landscape of data proliferated through the internet and in other published textual sources, such studies could also shed light on what kind of text is being read by the literate population, which styles are dominating, and what kind of language structures and content is being prioritised through the usage of these generative AI.  In addition, studies such as these can inform readers of potential heuristics to use when reading a text, influencing their intuition to scrutinise a text and potentially infer their origin from a machine. Analysing the characteristics of LLM-generated data enhances the ability to shed light on the mechanisms of LLMs, potentially contributing to transparency and LLM-probing research within the NLP space. 

The study of the linguistic properties of LLM-generated text has emerged from this recent phenomenon as a growing topic of natural language processing (NLP).  Recent research has focused primarily on English-language textual analysis as the dominant type of data analysed. However, there has been thus far limited but emerging research into the characteristics of LLM-generated text in languages other than English, such as Swedish. Hence, there is a lot of opportunity for exploration of this topic in many other languages. 

In addition, there has been emerging research into detection methods of LLM-generated text, both by AI detection tools and by humans. These studies help with development of AI detection tools and to further educate the general public in building their intuition of LLM-generated text. There is recent evidence that LLM-generated text has changed the make-up of text on the internet (INSERT SOURCES), with an estimated (SOURCE) \% of the internet estimated to be made up of LLM-generated text (although this figure is only estimated as it is difficult to gauge). AI-detection tools are useful in several contexts, including plagiarism detection and the propagation of LLM-generated text containing political propaganda and hate speech to influence users on the internet. These topics are pressing matters due to the number of users on the internet and the risk of political order being compromised and bias against groups being reinforced with the help of LLM tools, which can create massive amounts of data very quickly. 


\section{Research questions}

The intended purpose of this study is to investigate systematic differences (if any) between text generated
by large language models (LLMs) and text written by humans in the Swedish language. In addition, the
study will investigate whether there are any significant differences in the ability to detect LLM-generated
Swedish text between humans and LLM tools. The study will include both a formal and informal text
component, that is, linguistic analyses and AI-detection experiments will be performed on two sources
of data from formal and informal domains.

\begin{enumerate}
\item Are there systematic differences between LLM-generated and human-written formal text in
Swedish?
\item Are there systematic differences between LLM-generated and human-written informal text in
Swedish?
\item Are there differences between humans and AI-based detection tools in their ability to detect LLM-
generated Swedish text?
\end{enumerate}

\section{Experimental-set up of the thesis}

The thesis is comprised of two experiments: 1) textual analysis, and 2) LLM detection. The first experiment focuses on how human-written and LLM-generated text differ in surface-level and content-level characteristics. There are two domains of text used: 1) a formal domain comprised of Bachelor's level thesis abstracts, and 2) an informal domain of internet thread comments. The LLM used to generate the texts is OpenAI's GPT-5.2. 
The second experiment focuses on the differences in ability and thinking process between humans and LLMs in detecting LLM-generated text. The LLM of focus is DeepSeek (? plus model version), and the accuracy and reasoning process will be evaluated for both Swedish-native speakers and this LLM in detecting both the formal and informal data.



\chapter{Background}

This chapter will provide an overview of previous studies related to this topic. First, it will outline the development of large language models and their text generation abilities for English and Swedish. Second, it will describe previous research comparing AI-generated text with human-written text. The chapter will then outline tools and human participant studies related to AI-generated text detection. 
\section{Large Language Models and text generation}

\subsection{Large Language Model architectural development}


\paragraph{}

At the time of writing this thesis, the dominant underlying architecture of LLMs is the transformer architecture, first described in \textcite{vaswani2017attention}. In this study, the authors designed a basic neural network architecture with a multi-head attention mechanism and no recurrent unit structure. The transformer architecture has demonstrated superior performance and decreased training time compared to its recurrent neural network predecessors. The strength of transformers lies in their ability to handle long dependencies between different elements of input sequences, meaning that it can calculate how relevant different parts of input are to one another, no matter how far apart one element is from another within a sequence. As a result, transformer networks can be adapted to a variety of tasks, including optical recognition, speech processing, and text generation \parencite{islam2023}. 
\paragraph{}
Specifically look at LLM architecture and text generation itself, its relationship with training on data.

\subsection{LLMs for the English language}

English has a vast amount of data - list major datasets used for LLM development in English.

Talk about compositions of English-dominant LLMs - data type

\subsection{LLMs for the Swedish language}

Ohman et al. (2023) curated a 1.2 terabyte dataset of text in all Scandinavian languages in addition to high-quality English data, called the Nordic Pile dataset. Holmstrom et al. (2023) initiated an investigation of the extent to which pre-training GPT-3 LLMs on Swedish only is advantageous over leveraging pre-trained knowledge (from English). It also examined whether English-based pre-trained models demonstrated unwanted transfer of Anglophone cultural aspects over to the Swedish context.  The findings of Holmstrom et al. (2023) demonstrated that LLMs exposed to limited Swedish during training can be highly capable and transfer competencies from English "off-the-shelf", including mathematical reasoning and Swedish cultural adaptation. In addition, the study showed that GPT-3 models did not appear to exhibit unwanted over-generalisation of cultural aspects to Swedish. 

\section{Textual comparison of LLM-generated and human-written text}

Several studies have already done that analyse the linguistic characteristics of LLM-generated text in English. 

\subsection{Linguistic properties of LLM-generated text}

Several studies have already done that analyse the linguistic characteristics of LLM-generated text in English. 
\paragraph{}
Munoz-Ortiz et al. (2024) compared the linguistic components of human and LLM-generated news articles and compared their patterns among six generative LLMs. For their prompting strategy, they asked the models to generate a news paragraph based on the news article headline its first words. For their linguistic analysis, they compared the human-written and LLM-generated text groups through differences in their linguistic structures, including vocabulary distribution, part-of-speech tag distributions, dependency arc lengths, constituent lengths, cosine similarity, and emotional category tagging. The results of this analysis showed significant differences between the human-written and LLM-generated text on all textual analysis metrics. At the lexical level, 5 of the 6 LLM models outputted text with a more restricted vocabulary than the human-written text. For semantics, the models were found to have less of an inclination to display aggressive negative emotions (such as fear or anger) than humans. 
\paragraph{}
Liao et al. (2023) analysed the differences between human-written and ChatGPT-generated text within the medical domain. Medical texts were compared by their vocabulary distribution, part-of-speech tag distribution, dependency arcs, sentiment analysis, and perplexity. In addition, they implemented machine learning methods to detect ChatGPT-generated medical text. The results of the study found that human-written medical texts contained more concrete information, had higher vocabulary diversity, and contained information that human experts would consider more useful. The ChatGPT-generated medical texts appeared more fluent and logical, but expressed more generalised terminology rather than specific, effective information for the problem's context. Overall, this would point to the human-written medical text being more tailored to the situational modelling of the medical case (situational modelling defined in  Mahowald et al., 2024) rather than the LLM-generated text.

\paragraph{}
- discuss differences in text stylistic 

\subsection{Other properties of LLM-generated text}
Previous work has also shown evidence of a distinction between the reasoning patterns of LLM tools when compared to those of humans. (CITE) evaluated the medical reasoning capabilities of DeepSeek R1, computing its outputs to the reasoning patterns of medical domain human experts. The study performed qualitative and quantitative analyses of AI- and human-medical reasoning for 100 clinical cases from the MedQA dataset. The study found a 93\% diagnostic accuracy for DeepSeek R1's performance. Many of the LLM's errors were attributed to over-focus on the details of a medical problem, rather than a more holistic and intuitive (?) understanding by humans, as well as an insufficient understanding of medical protocols when compared to human experts. The study also found a strong correlation between the length of output and the error rate of LLMs; t hat is, results from LLMs that were longer were more likely to contain incorrect diagnoses or medical details.

- discuss differences in text stylistic 


\section{LLM-generated text detection studies}

\subsection{Human participant AI-detection studies}
Landberg (2024) investigated the effectiveness of AI chatbots to detect texts created by AI in Swedish and English. The main component of the study  A recent study investigated LLM reasoning in Swedish. Moell and Boye (2025) studied the performance of LLMs on language complexity measurement tasks (LIX readability metric and Average Dependency Distance). The results of the study found that all models demonstrated at least come capacity for each of these tasks, but that LLMs remain prone to errors in explicit calculations and syntactic analysis. 
\subsection{LLM-tool detection by LLMs studies}


\section{Background summary}

- add a breakdown of gpt-5.2 Swedish training data

\chapter{Methodology}

The methodology is comprised of two main parts: 1) a comparison of textual elements between human-written Swedish and LLM-generated Swedish, and 2) a comparison between the detection rates of humans and LLM-detection of LLM-generated Swedish. Both parts of the methodology comprise of a human component and an LLM-generated component.



\section{Data collection}

The data collection of this thesis formed the data source of both the textual analysis and LLM detection experiments. The dataset comprises of  1) a set of texts in a formal register, collected through abstracts; and 2) a set of texts in an informal register, collected through forum comments. These two registers are selected to measure differences in textual elements and detection rates of LLM-generated text within two different contexts.

\subsection{Human-written formal data}

The formal human-written dataset comprises of (n=142(?)) Bachelor's level thesis titles, abstracts and keywords sourced from the Digitala Vetenskapliga Arkivet (DiVA) portal. All Swedish abstracts and keywords were written originally in Swedish. All abstracts were sourced between 2007-2017, to negate the chance of of any transformer-based generative models influencing any part of the dataset's titles or abstracts. The topics of the thesis abstracts were sampled randomly across (insert number) topics. On average, the thesis abstracts were  (n=words) in length. 

\subsection{Human-written informal data}

-- THE DATA COLLECTION FOR THIS NEEDS TO BE RE-DONE

The informal human-written dataset comprises of (34?) Swedish-language Reddit questions, with (n=insert) replies to each question, all scraped from r/Sweden and r/svenskpolitikpushctic Shift APIs. To mitigate the possibility of transformer-based generative AI  tools being involved in the creation of the subreddit's questions and replies, only questions and comments with a timestamp before January 2017 (2017-01-01) were scraped. To ensure that the post's titles were in question form and in Swedish, the scraper script checked for the presence of a question mark ("?") and used the <langdetect> Python library to identify the post as being written in Swedish ("sv"). On average, the Reddit comments 

\subsection{LLM-generated formal data}

The formal human-written dataset comprises of an LLM-generated version of each human-written sample of the DiVA Bachelor's-level thesis abstracts. The model used was OpenAI's GPT 5.2 model (<gpt-5.2>). The prompt used as input to generate the comments was:
\paragraph{}
"Titta på titeln och nyckelorden och skapa en sammanfattning i kandidatuppsatsstil med dina egna ord: <title, keywords>."
\paragraph{}
English translation: Look at the title and keywords and create a summary of the thesis abstract in your own words: <title, keywords>



\subsection{LLM-generated informal data}


\section{Experiment 1: textual analysis}

The aim of the textual analysis experiments is to investigate meaningful differences between LLM-generated and human-written Swedish. This method is based on the hypothesis that there are indeed meaningful differences in the writing style between the two textual sources in Swedish. The method aims to cover both surface-level analyses of the text, that is, investigating structural differences between the word forms and surface-level composition between the two modes; as well as content-level analysis, that is, sentiment analysis, entities, and themes of the texts. 
\subsection{Surface-level analysis}
For the surface-level analysis, the spaCy package was used (\texttt{sv\_core\_news\_sm}) to collect distributions of: vocabulary distribution, part-of-speech distribution, dependency analysis (sprakbanken?), dependency arch length, average word length, average sentence length. 

\subsection{Content-level analysis}

For the content analysis, the study used the sentiment analysis, named entity recognition, and thematic clustering(?) techniques to uncover recurring ideas and themes within both the abstracts and comments. Hypothesis: this is going to be more important for the Reddit comments than the abstracts.


\section{Experiment 2: LLM-generated text detection}


\subsection{LLM-detection experimental set-up}

Another LLM (DeepSeek or whatever) to do AI detection  - quantitative measure (accuracy),  qualitative measure (reasoning) 

\subsection{Human detection} 
For the human detection experiment, a web game was created by the author with 15 Reddit-style forum threads with 4 comment replies per question, and 15 thesis abstract questions. The URL of this web game was shared in Swedish-speaking forums, with the only prerequisite screening for users of the game was native-level knowledge of Swedish.
\paragraph{}
Before the game started, the user was prompted to fill in information about 1) their age, 2) their occupation, and 3) the estimated number of hours of their LLM tool usage per week. Each thread and abstract was shown one at a time, with instructions given to the user to determine whether each comments or abstract was LLM-generated or human-written. The proportion of AI comments or abstracts to human-written was sampled at random. After each question was completed, feedback was given immediately to the user as to whether their response was correct or not before moving on to the next round. This immediate feedback characteristic was decided in order to increase the chance of engagement by the user until the very end. Each user was given the option to exit the game at any time. At the end of the game session (or upon exit), the user was shown their accuracy results and prompted to fill in an optional short reflection to explain their reasoning process for LLM-detection throughout the game.
\paragraph{}
There were (n=?) game sessions submitted through the web game, each of which was stored in a Google Sheet form with information about their usage time-stamp, question identification number, which questions were determined as human, and accuracy scores. The results were measured in terms of accuracy (F-score too?), and qualitative analysis of any reasoning descriptions sent through the game. 

\subsection{AI detection}



\chapter{Results}
text here


\chapter{Discussion}
text here


\chapter{Conclusion}
text here

%\nocite{*}

\printbibliography
% at start yields errors because there are no citations and the bibliography file is empty

\end{document}
